% ============================================================================ %
%
%           Šablona bakalářské/diplomové práce UTB ve Zlíně
%
% Autor:    2006/05, základ, Jozef Říha
%           2024/05, úpravy citování, Roman Šír (r_sir@utb.cz)
%           2014-2025, Ing. Pavel Tomášek, Ph.D. (tomasek@utb.cz)
%
% Verze:    2025-09-16
%
% Kódování: UTF-8 (kontrolní řetězec: žluťoučký kůň úpěl ďábelšké ódy)
%
% Sazba:    Doporučená kombinace TeXstudio a LaTeX (MikTeX), případně Overleaf.
%           Ukázka překladu z terminálu, či z příkazové řádky pro korektní vložení referencí:
%           `pdflatex thesis.tex && bibtex thesis && pdflatex thesis.tex && pdflatex thesis.tex`
%
% Tip:      Ve správně vysázeném českém textu by na konci řádku neměla zůstat 
%           samotná jednopísmenná předložka či spojka. Na takové místo se vkládá
%           nezalomitelná mezera pomocí symbolu ~.
%
% Pozor:    Vzhledem k požadovanému standardu PDF/A nesmí vložené obrázky
%           obsahovat alfa kanál (průhlednost).
%
% ============================================================================ %

% Definice vzhledu a nastavení se nachází přímo v následující speciální třídě
\documentclass{template-definitions/thesistbuinzlin}

% Uživatelské definice -- upravte dle potřeby
\setthesislanguage{CZ}
% CZ   -- práce bude v českém jazyce
% EN   -- práce bude v anglickém jazyce

\setfaculty{FAI}
  % FAI  -- pro Fakultu aplikované informatiky
  % FAME -- pro Fakultu managementu a ekonomiky
  % FHS  -- pro Fakultu humanitních studií
  % FLKR -- pro Fakultu logistiky a krizového řízení
  % FMK  -- pro Fakutlu mutimediálních komunikací
  % FT   -- pro Fakultu technologickou
  % UNI  -- pro Univerzitní institut

\setthesistype{DP}
  % BP   -- bakalářská práce
  % DP   -- diplomová práce

\setthesisyear{2026}
  % zadejte rok namísto "xxxx"

\setcitationmode{numbers}
  % authoryear -- odkazování na zdroje bude v módu (autor, rok), tzv. harvardský systém, vhodné v případě fakult FAME, FHS, FLKŘ, FMK, FT
  % numbers    -- odkazování na zdroje bude v číselném módu [číslo], vhodné v případě fakult FAI, FT

% Lze přidat vertikální odsazení nad (prvni parametr) a pod (druhý parametr) obrázky, tabulky i rovnice/soustavy rovnic
\setmarginaroundfigures{0mm}{0mm}
\setmarginaroundtables{0mm}{0mm}
\setmarginaroundequations{0mm}{0mm}

\setauthor{Bc. Josef Susík}

\settitlecz{Digitalizace domácí knihovny}
\settitleen{Thesis Title in English (max. 7 rows)}  % Needed only when writing in English

\setabstractcz{Text abstraktu v jazyce práce.}

\setabstracten{Text abstraktu ve světovém jazyce (angličtině).}

\setkeywordscz{klíčové slovo, klíčové slovo}

\setkeywordsen{keyword, keyword}

\setuppdfa

%\usecolouredlinks  % Volitelné, může být nápomocné při tvorbě a kontrole výsledného dokumentu

% ============================================================================ %

\begin{document}

\maketitle

\insertassignment{graphics/placeholder-scan.pdf}

\insertdeclaration{CZ}
  % CZ   -- doporučované, i pokud práci píše český/slovenský student v angličtině
  % EN   -- variant only for international students

\insertabstractandkeywords

% ============================================================================ %

\cleardoublepage
Zde je místo pro případné poděkování, popř. motto, úryvky knih atp.
Pokud poděkování nechcete vložit, smažte celý tento blok 
(mezi příkazy cleardoublepage, včetně těchto příkazů).

{\bigskip
\color{red}UPOZORNĚNÍ: důležitou součástí této šablony je manuál s návodem. Důrazně proto doporučujeme přečíst si manuál k šabloně.}
\cleardoublepage

% ============================================================================ %

\insertcontents  % Obsah je generován automaticky

% ============================================================================ %

\insertlistoffigures  % Seznam je generován automaticky

\insertlistoftables  % Seznam je generován automaticky

%\insertlistoflisting  % Seznam je generován automaticky, vhodné jen pokud máte v dokumentu fragmenty kódu

\input{text/abbreviations.tex}  % Seznam zkratek

\insertlistofappendices  % Seznam příloh je generován automaticky

% ============================================================================ %

\paragraphindentationstart  % Nastaví odsazování odstavců dle zvoleného jazyka
\cleardoublepage					  % Další text (úvod) až na liché straně

% Text práce
% ============================================================================ %
% Encoding: UTF-8 (žluťoučký kůň úpěl ďábelšké ódy)
% ============================================================================ %

\sectionwithoutnumbering{Úvod}



% ============================================================================ %


% Pouze u vybraných fakult (např. FAME):
%   vložení samostatné nečíslované kapitoly týkající se cílů a metod
\iffame \input{text/text-chapter-00-objectives.tex} \fi

\part{TEORETICKÁ ČÁST}
% ============================================================================ %
% Encoding: UTF-8 (žluťoučký kůň úpěl ďábelšké ódy)
% ============================================================================ %

% ============================================================================ %
\section{Digitální knihovny}
Text \dots


\subsection{Definice digitální knihovny}
Text \dots


\subsection{Historický vývoj digitálních knihoven}
Text \dots


\subsection{Historický vývoj digitálních knihoven}
Text \dots


\subsection{Metadata a katalogizační standardy}
Text \dots


\subsection{Současné komerční a open-source platformy}
Text \dots


% ============================================================================ %

% ============================================================================ %
% Encoding: UTF-8 (žluťoučký kůň úpěl ďábelšké ódy)
% ============================================================================ %

% ============================================================================ %
\section{Doporučovací systémy}
Text \dots


\subsection{Doporučovací systémy}
Text \dots


\subsection{Typy doporučování}
Text \dots


\subsection{Algoritmy používané v doporučování}
Text \dots


\subsection{Algoritmy používané v doporučování}
Text \dots


\subsection{Personalizace obsahu v digitálních knihovnách}
Text \dots


% ============================================================================ %


\part{PRAKTICKÁ ČÁST}
\input{text/text-chapter-03.tex}

\input{text/text-chapter-99-conclusion.tex}

\paragraphindentationstop

% ============================================================================ %

% Využívají se bibliografické údaje uložené v souboru tex/references.bib.
% Ty se automaticky upravuji dle zvolené citační normy (ČSN ISO 690:2022).
\insertlistofreferences

\input{text/appendices.tex}  % Vložení příloh

% ============================================================================ %

\end{document}

% ============================================================================ %
